\section{Geometry \texorpdfstring{\&}{&} Misc}
\subsection{Geometry Formulas}
Triangle:

Side length $a, b, c$

Area: $\frac{1}{2} \times base \times height$

$s = \frac{a+b+c}{2}$, Area $= \sqrt{s(s-a)(s-b)(s-c)}$

Inradius: $r = \frac{Area}{s}$

Circumradius: $R = \frac{abc}{4 \times Area}$

Circles:

Arc Length: $L = \theta r$

Sector Area: $Area = \frac{1}{2} r^2 \theta$

Segment Area (Area of the part cut off by a chord): $Area = \frac{1}{2} r^2 (\theta - \sin\theta)$

\begin{itemize}
  \item \textbf{Rhombus:}
  \begin{itemize}
    \item Area: $\frac{1}{2} \times d_1 \times d_2$ (where $d_1$ and $d_2$ are diagonals).
  \end{itemize}
  \item \textbf{Trapezium:}
  \begin{itemize}
    \item Area: $\frac{1}{2} \times (\text{Base}_1 + \text{Base}_2) \times \text{Height}$.
  \end{itemize}
\end{itemize}
3. \textbf{Regular Polygon} (\textit{n sides}):
\begin{itemize}
  \item Perimeter: $n \times a$
  \item Area: $\frac{1}{4}na^2 \cot\!\left(\frac{\pi}{n}\right)$ (where $a$ is the side length).
\end{itemize}

\textbf{Curved Shapes}

1. \textbf{Circle}
\begin{itemize}
  \item Circumference: $2\pi r$
  \item Area: $\pi r^2$
\end{itemize}
2. \textbf{Ellipse}
\begin{itemize}
  \item Area: $\pi a b$ (where $a$ is the semi-major axis and $b$ is the semi-minor axis).
\end{itemize}

\textbf{$\checkmark$ Polygon}

\begin{itemize}
  \item Interior angle sum:
  $(n-2) \times 180^{\circ}$

  \item Each internal angle (regular):
  $\frac{(n-2)\times180^{\circ}}{n}$
\end{itemize}

4. \textbf{Tetrahedron}
\begin{itemize}
  \item Volume: $\frac{1}{6} \times \text{Base Area} \times \text{Height}$.
\end{itemize}

\textbf{Curved 3D Shapes}

1. \textbf{Sphere}
\begin{itemize}
  \item Surface Area: $4\pi r^2$
  \item Volume: $\frac{4}{3}\pi r^3$
\end{itemize}
2. \textbf{Cylinder}
\begin{itemize}
  \item Surface Area: $2\pi r(r+h)$
  \item Volume: $\pi r^2 h$
\end{itemize}
3. \textbf{Cone}
\begin{itemize}
  \item Surface Area: $\pi r(r+l)$, where $l=\sqrt{r^2+h^2}$
  \item Volume: $\frac{1}{3}\pi r^2 h$.
\end{itemize}

\textbf{Other Shapes}

1. \textbf{Torus}
\begin{itemize}
  \item Surface Area: $4\pi^2 Rr$
  \item Volume: $2\pi^2 Rr^2$, where $R$ is the distance from the center of the tube to the center of the torus, and $r$ is the radius of the tube.
\end{itemize}
\subsection{Closest Pair of Points}
\begin{minted}{c++}
struct Res {
  int64_t d;
  P a, b;
};
int64_t sq(int64_t x) { return x * x; }
int64_t dist(P &a, P &b) {
  return sq(a.x - b.x) + sq(a.y - b.y);
}
// p = py = tmp, all are same sorted by x co-ordi.;
Res closest_pair(vector<P>& p, vector<P>& py,
vector<P>& tmp, int l, int r) {
  int n = r - l + 1;
  if (n <= 1) return {LLONG_MAX, {}, {}};
  if (n == 2) {
    if (py[l].y > py[l + 1].y)
      swap(py[l], py[l + 1]);
    return {dist(p[l], p[l + 1]), p[l], p[l + 1]};
  }
  int m = (l + r) >> 1;
  int midx = p[m].x;
  Res res_l = closest_pair(p, py, tmp, l, m);
  Res res_r = closest_pair(p, py, tmp, m + 1, r);
  Res res = (res_l.d < res_r.d ? res_l : res_r);
  for (int i = l, j = m + 1, k = l; k <= r; k++) {
    if (j > r || (i <= m && py[i].y < py[j].y))
      tmp[k] = py[i++];
    else tmp[k] = py[j++];
  }
  for (int i = l; i <= r; i++) py[i] = tmp[i];
  vector<P> stripe;
  for (int i = l; i <= r; i++) {
    if (sq(py[i].x - midx) < res.d)
      stripe.push_back(py[i]);
  }
  for (int i = 0; i < (int)stripe.size(); i++) {
    for (int j = i + 1; j < (int)stripe.size(); j++) {
      if (sq(stripe[i].y - stripe[j].y) >= res.d)
        break;
      int64_t d = dist(stripe[i], stripe[j]);
      if (d < res.d) res = {d, stripe[i], stripe[j]};
    }
  }
  return res;
}
\end{minted}
\begin{minted}{c++}
long double L, W;
if(!(cin >> L >> W)) return 0;
// hypotenuse of each right triangle
long double h = sqrtl(L*L + W*W);
// perimeter for incenter-weights: P = L + W + h
long double P = L + W + h;
// incenter of triangle with vertices
// (0,0),(L,0),(L,W)
long double x1 = L*(L + h)/P;
long double y1 = L*W/P;
// incenter of triangle with vertices
// (0,0),(0,W),(L,W)
long double x2 = L*W/P;
long double y2 = W*(W + h)/P;
long double dx = x1 - x2;
long double dy = y1 - y2;
long double D = sqrtl(dx*dx + dy*dy);
cout << fixed << setprecision(6) << (double)D
<< "\n";
\end{minted}
\begin{minted}{c++}
long double r1, r2, r3;
cin >> r1 >> r2 >> r3;
// side-lengths of triangle ABC
long double a = r2 + r3;
long double b = r3 + r1;
long double c = r1 + r2;
long double s = (a + b + c) / 2.0L;
long double T = sqrtl( s*(s - a) * (s - b) * (s - c) );
// helper to compute angle opposite side a
auto angle = [&](long double opp, long double
x, long double y){
  long double cosv = (x*x + y*y - opp*opp) /
  (2.0L * x * y);
  // guard numerical drift
  cosv = min((long double)1.0, max((long
  double)-1.0, cosv));
  return acosl(cosv);
};
long double A = angle(a, b, c);
long double B = angle(b, c, a);
long double C = angle(c, a, b);
// subtract the sectors
long double sectorSum = 0.5L*(r1*r1*A +
r2*r2*B + r3*r3*C);
long double blackArea = T - sectorSum;
cout << fixed << setprecision(6) <<
(double)blackArea << "\n";
return 0;
}
\end{minted}
\subsection{Line Intersection}
\begin{minted}{c++}
// Intersecting point of two lines (a1-a2, b1-b2)
using P = pair<ll,ll>;
auto cross = [&](P a, P b){ return
(__int128)a.first*b.second -
(__int128)a.second*b.first; };
auto intersect = [&](P a1, P a2, P b1, P b2) ->
optional<pair<long double,long double>> {
  P da = {a2.first - a1.first,
  a2.second - a1.second};
  P db = {b2.first - b1.first,
  b2.second - b1.second};
  P dc = {b1.first - a1.first,
  b1.second - a1.second};
  __int128 D = cross(da, db);
  if (!D) return {}; // parallel or coincident
  long double t = (long double)cross(dc, db) /
  (long double)D;
  return {{a1.first + da.first * t,
  a1.second + da.second * t}};
};
// Example: // if (auto p = intersect(a1,a2,b1,b2))
// cout << llround(p->first) << " " <<
// llround(p->second) << "\n";
// else
// cout << "No intersection\n";
// For segments add this before returning
long double u = (long double)cross(dc, da) / (long
double)D;
if (t < 0 || t > 1 || u < 0 || u > 1) return {};
// no segment intersection
\end{minted}
\subsection{Polygon Operations}
\begin{minted}{c++}
const ld EPS = 1e-9;
const ld PI = acos(-1.0);
int sgn(ld x) { return (x > EPS) - (x < -EPS); }
struct P {
  using T = int64_t;
  T x, y;
  void read() { cin >> x >> y; }
  P operator - (const P &b) const {
    return P{x - b.x, y - b.y}; }
  P operator + (const P &b) const {
    return P{x + b.x, y + b.y}; }
  void operator -= (const P &b) {
    x -= b.x; y -= b.y; }
  void operator += (const P &b) {
    x += b.x; y += b.y; }
  bool operator < (const P &b) const {
    return x == b.x ? y < b.y : x < b.x; }
  bool operator == (const P &b) const {
    return x == b.x && y == b.y; }
  T operator * (const P &b) const {
    return x * b.y - y * b.x; }
  T operator ^ (const P &b) const {
    return x * b.x + y * b.y; }
  T dot (P &a, P &b) const {
    return (a - *this) ^ (b - *this); }
  T cross (P &a, P &b) {
    return (a - *this) * (b - *this); }
  long double slope() {
    return (long double)y / x; }
  ld sq() const { return x * x + y * y; }
  ld abs() const { return sqrtl(sq()); }
  P rot() const { return P(-y, x); } // Rotate 90 deg CCW
  P rot(ld a) const {
    return P(x*cos(a) - y*sin(a),
    x*sin(a) + y*cos(a)); }
};
\end{minted}
\begin{minted}{c++}
bool collinear(P &a, P &b, P &c) {
  return (b - a) * (c - a) == 0;
}
inline bool on_segment(P &p, P &q, P &r) {
  if (!collinear(p, q, r)) return false;
  return r.x <= max(p.x, q.x) && r.x >= min(p.x,
  q.x) && r.y <= max(p.y, q.y) &&
  r.y >= min(p.y, q.y);
}
ld dist(P a, P b) { return (a - b).abs(); }
// --- Lines & Segments ---
// Project C onto line AB
P project(P a, P b, P c) {
  return a + (b - a) * (((c - a) ^ (b - a)) /
  (b - a).sq());
}
// Reflect C across line AB
P reflect(P a, P b, P c) {
  return project(a, b, c) * 2 - c;
}
// Dist from C to infinite line AB
ld dist_pt_line(P a, P b, P c) {
  return fabsl((b - a) * (c - a)) / (b - a).abs();
}
// Dist from C to segment AB
ld dist_pt_seg(P a, P b, P c) {
  if (((c - a) ^ (b - a)) <= 0) return (c - a).abs();
  if (((c - b) ^ (a - b)) <= 0) return (c - b).abs();
  return dist_pt_line(a, b, c);
}
\end{minted}
\begin{minted}{c++}
// Check if P is on segment AB
bool on_seg(P a, P b, P p) {
  return !sgn((b - a) * (p - a)) &&
  ((p - a) ^ (p - b)) <= EPS;
}
// Line-Line Intersection (returns false if parallel)
bool line_int(P a, P b, P c, P d, P &res) {
  ld cp = (b - a) * (d - c);
  if (!sgn(cp)) return false;
  res = a + (b - a) * (((d - c) * (c - a)) / cp);
  return true;
}
// Segment-Segment Intersection
bool seg_int(P a, P b, P c, P d) {
  auto cp = [&](P x, P y, P z) {
    return sgn((y - x) * (z - x)); };
  if (on_seg(a, b, c) || on_seg(a, b, d) ||
  on_seg(c, d, a) || on_seg(c, d, b)) return true;
  return cp(a, b, c) != cp(a, b, d) &&
  cp(c, d, a) != cp(c, d, b);
}
long double polygon_area(vector<P> &p) {
  int64_t area = 0, n = p.size();
  for (int i = 0; i < n; ++i)
    area += p[i] * p[(i + 1) % n];
  return abs(area) / 2;
}
\end{minted}
\begin{minted}{c++}
vector<P> convexHull(vector<P> p) {
  int n = p.size(), k = 0;
  if (n <= 2) return p;
  vector<P> h(2 * n);
  sort(p.begin(), p.end());
  for (int i = 0; i < n; h[k++] = p[i++])
    while (k >= 2 && h[k - 2].cross(h[k - 1],
    p[i]) <= 0) k--;
  for (int i = n - 2, t = k + 1; i >= 0;
  h[k++] = p[i--])
    while (k >= t && h[k - 2].cross(h[k - 1],
    p[i]) <= 0) k--;
  h.resize(k - 1);
  return h;
}
int is_point_in_polygon(const vector<P>& p, P
q) {
  int c = 0;
  for (int i = 0, n = p.size(); i < n; ++i) {
    P a = p[i], b = p[(i + 1) % n];
    if ((b - a) * (q - a) == 0 && min(a.x, b.x) <= q.x
    && q.x <= max(a.x, b.x) && min(a.y, b.y) <= q.y
    && q.y <= max(a.y, b.y)) return 2;
    if ((a.y > q.y) != (b.y > q.y) &&
    ((b - a) * (q - a) > 0) == (a.y < b.y))
      c = !c;
  }
  return c;
}
double angle_at_vertex(P &a, P &b, P &c) {
  double ang_rad = atan2l(fabs(b.cross(a, c)),
  b.dot(a, c));
  return ang_rad * 180.0L / acosl(-1.0L);
}
ld dist_point_line(P a, P b, P c) {
  ld area2 = fabsl((b - a) * (c - a));
  ld base = hypotl(b.x - a.x, b.y - a.y);
  return area2 / base;
}
\end{minted}
\subsection{Mo's Algorithm}
\begin{minted}{c++}
sort(cu.begin(), cu.end(), [&](auto &x, auto &y){
  if (x.l / B != y.l / B) return x.l / B < y.l / B;
  return (x.l / B) % 2 ? x.r > y.r : x.r < y.r;
});
int R = -1, L = 0;
while (R < r) add(++R);
while (L > l) add(--L);
while (R > r) rem(R--);
while (L < l) rem(L++);
\end{minted}
\subsection{Fast Factorization (Pollard Rho)}
\begin{minted}{c++}
#define ull uint64_t
struct Montgomery {
  ull n, nr;
  Montgomery(ull n) {
    this->n = n, nr = 1;
    for (int i = 0; i < 6; i++)
      nr *= 2 - n * nr;
  }
  ull reduce(__uint128_t x) const {
    ull q = ull(x) * nr,
    m = ((__uint128_t)q * n) >> 64;
    return (x >> 64) + n - m;
  }
  ull multiply(ull x, ull y) {
    return reduce((__uint128_t)x * y); }
};
ull f(ull x, Montgomery& m, ull a) {
  return m.multiply(x, x) + a; }
ull diff(ull a, ull b) {
  return a > b ? a - b : b - a; }
const int M = 1024;
const ull SEED = 42;
ull find_factor(ull n, ull x0 = 2, ull a = 1)
{
  Montgomery m(n);
  ull x = SEED;
  for (int l = M; l < (1 << 20); l *= 2) {
    ull y = x, p = 1;
    for (int i = 0; i < l; i += M) {
      for (int j = 0; j < M; j++)
        x = f(x, m, a), p = m.multiply(p,
        diff(x, y));
      if (ull g = gcd(p, n); g != 1)
        return g;
    }
  }
  return n;
}
int32_t main() {
  ll n;
  cin >> n;
  ll factor = find_factor(n);
  if (factor != 1) {
    cout << "Factor found: " << factor << endl;
  } else {
    cout << "No factor found" << endl;
  }
  return 0;
}
\end{minted}
\subsection{Point and Vector}
\begin{minted}{c++}
const double EPS = 1e-9;
struct Point {
  double x, y;
  Point operator+(const Point& other) const {
    return {x + other.x, y + other.y}; }
  Point operator-(const Point& other) const {
    return {x - other.x, y - other.y}; }
  Point operator*(double s) const { return {x * s,
  y * s}; }
  Point operator/(double s) const { return {x / s,
  y / s}; }
  bool operator<(const Point& other) const {
    if (abs(x - other.x) > EPS) return x < other.x;
    return y < other.y - EPS;
  }
  bool operator==(const Point& other) const {
    return abs(x - other.x) < EPS && abs(y -
    other.y) < EPS;
  }
};
typedef Point Vector;
double dot(Vector a, Vector b) { return a.x * b.x +
a.y * b.y; }
double norm_sq(Vector a) { return dot(a, a); }
double length(Vector a) { return
sqrt(norm_sq(a)); }
double cross(Vector a, Vector b) { return a.x * b.y
- a.y * b.x; }
\end{minted}
\begin{minted}{c++}
// Returns: > 0 for Left Turn, < 0 for Right Turn,
// 0 for Collinear
double orientation(Point p, Point a, Point b) {
  return cross(a - p, b - p);
}
double distToLine(Point p, Point a, Point b) {
  return abs(cross(p - a, b - a)) / length(b - a);
}
double distToSegment(Point p, Point a, Point b) {
  if (dot(p - a, b - a) < 0) return length(p - a);
  if (dot(p - b, a - b) < 0) return length(p - b);
  return distToLine(p, a, b);
}
double polygonArea(const vector<Point>& p) {
  double area = 0;
  for (int i = 0; i < (int)p.size(); i++) {
    area += cross(p[i], p[(i + 1) % p.size()]);
  }
  return abs(area) / 2.0;
}
\end{minted}
\begin{minted}{c++}
void convex_hull(vector<pt>& a, bool
include_collinear = false) {
  pt p0 = *min_element(a.begin(), a.end(), [](pt
  a, pt b) {
    return make_pair(a.y, a.x) < make_pair(b.y,
    b.x);
  });
  sort(a.begin(), a.end(), [&p0](const pt& a, const
  pt& b) {
    int o = orientation(p0, a, b);
    if (o == 0)
      return (p0.x-a.x)*(p0.x-a.x) +
      (p0.y-a.y)*(p0.y-a.y)
      < (p0.x-b.x)*(p0.x-b.x) +
      (p0.y-b.y)*(p0.y-b.y);
    return o < 0;
  });
  if (include_collinear) {
    int i = (int)a.size()-1;
    while (i >= 0 && collinear(p0, a[i], a.back()))
      i--;
    reverse(a.begin()+i+1, a.end());
  }
  vector<pt> st;
  for (int i = 0; i < (int)a.size(); i++) {
    while (st.size() > 1 && !cw(st[st.size()-2],
    st.back(), a[i], include_collinear))
      st.pop_back();
    st.push_back(a[i]);
  }
  if (include_collinear == false && st.size() == 2
  && st[0] == st[1])
    st.pop_back();
  a = st;
}
\end{minted}
